\section*{Hoofdstuk \ref{ch:b3:electrode-kinetics} --- Elektrodekinetiek en elektroanalyse}

\begin{solution}{exo:b3:electrode-kinetics:1}
Kathodische helling $2.303RT/\alpha F = \qty{118}{mV}$ geeft $\alpha = 59.2/118 = 0.50$; $j_0 = \qty{2.0e-6}{A.cm^{-2}}$,
de afsnijding bij $\eta = 0$.
\end{solution}

\begin{solution}{exo:b3:electrode-kinetics:2}
$i_0 = \qty{2.0e-4}{A}$; $R_{\mathrm{ct}} = RT/(Fi_0) = 0.02569/\num{2.0e-4} = \qty{128}{\ohm}$.
\end{solution}

\begin{solution}{exo:b3:electrode-kinetics:3}
Voor \ce{MgSO4} is $I = \frac12(4c + 4c) = 4c$. Bij \qty{1.0}{mmol/L} is $I = \qty{4.0}{mol.m^{-3}}$ en $\kappa^{-1} =
\qty{0.96}{nm} \times \sqrt{100/4} = \qty{4.8}{nm}$; bij \qty{0.10}{mol/L} is $I = \qty{400}{mol.m^{-3}}$, $\kappa^{-1} = \qty{0.48}{nm}$.
\end{solution}

\begin{solution}{exo:b3:electrode-kinetics:4}
$i_{\mathrm c} = C_{\mathrm{dl}}Av = \num{20e-6} \times 0.0707 \times 0.1 = \qty{0.14}{\micro A}$, minder dan 1\,\% van de piek. Bij
\qty{10}{V/s} is ze \qty{14}{\micro A}, terwijl de piek maar groeit tot $19\sqrt{100} = \qty{190}{\micro A}$:
de laadstroom groeit als $v$, de faradische als $\sqrt v$.
\end{solution}

\begin{solution}{exo:b3:electrode-kinetics:5}
$D = \pi\bigl(i\sqrt t/nFAc^*\bigr)^2 = \pi(\num{12.2e-6}/(96\,485 \times 0.0707 \times \num{1.00e-6}))^2 =
\qty{1.0e-5}{cm^2.s^{-1}}$.
\end{solution}

\begin{solution}{exo:b3:electrode-kinetics:6}
$FAc^* = \qty{1.364e-2}{C.cm^{-1}}$, $F/RT = \qty{38.92}{V^{-1}}$. Bij \qty{50}{mV/s}: $\sqrt{38.92 \times 0.05 \times \num{7.0e-6}} =
\num{3.69e-3}$, $i_{\mathrm p} = 0.4463 \times \num{1.364e-2} \times \num{3.69e-3} = \qty{22.5}{\micro A}$; bij \qty{500}{mV/s} $\sqrt{10}$ keer
meer, \qty{71}{\micro A}.
\end{solution}

\begin{solution}{exo:b3:electrode-kinetics:7}
(a) Reversibel. (b) \omterm{def:b3:electrode-kinetics:reversibility}{Quasi-reversibel}: de scheiding groeit met de
scansnelheid. (c) Een chemische stap volgt op de elektronenoverdracht
(EC): bij trage scans wordt het product verbruikt vóór de terugscan; een
snelle scan blijft de chemie voor.
\end{solution}

\begin{solution}{exo:b3:electrode-kinetics:8}
$K_{ij}a_j/a_i = \num{2e-4} \times 0.14/\num{4.0e-3} = 0.007$: de elektrode leest 0.7\,\% te hoog af.
\end{solution}

\begin{solution}{exo:b3:electrode-kinetics:9}
Kleinste kwadraten: $b = \qty{0.1093}{\micro A.L.\micro g^{-1}}$, $a = \qty{0.466}{\micro A}$, $s = \qty{0.011}{\micro A}$; $c_0 = a/b =
\qty{4.26}{\micro g.L^{-1}}$, met $u(c_0) = \frac sb\sqrt{\frac1n + \frac{\bar y^2}{b^2S_{xx}}} = \qty{0.12}{\micro g.L^{-1}}$
($n = 4$, $\bar y = 1.286$, $S_{xx} = 125$).
\end{solution}

\begin{solution}{exo:b3:electrode-kinetics:10}
$Q = 0.0500 \times 1800 = \qty{90.0}{C}$; $n(\ce{Cu}) = Q/2F = \qty{4.66e-4}{mol}$, \qty{29.6}{mg}. Het
resultaat hangt alleen af van de lading, gemeten uit een stroom en een
tijd, en van de constante van Faraday: er is geen standaard nodig, op
voorwaarde dat de elektrolyse volledig is en het stroomrendement
100\,\% bedraagt.
\end{solution}

\begin{solution}{exo:b3:electrode-kinetics:11}
$\sqrt{DFv/RT} = \sqrt{\num{1.0e-5} \times 38.92 \times 0.1} = \qty{6.24e-3}{cm.s^{-1}}$: $\Lambda = 1.6$, \omterm{def:b3:electrode-kinetics:reversibility}{quasi-reversibel}.
Bij \qty{10}{V/s} is $\Lambda = 0.16$, nog steeds \omterm{def:b3:electrode-kinetics:reversibility}{quasi-reversibel} maar
verder van de reversibele limiet; de golf zou pas onder ongeveer
\qty{1}{mV/s} reversibel zijn.
\end{solution}

\begin{solution}{exo:b3:electrode-kinetics:12}
$s/b = \qty{0.0838}{mM}$ en $b^2S_{xx} = \qty{203.5}{\micro A^2}$. Bij $\bar y$: $0.0838\sqrt{1 + 1/7} = \qty{0.090}{mM}$.
Bij \qty{18.0}{\micro A}: $0.0838\sqrt{1.143 + 83.0/203.5} = \qty{0.104}{mM}$. Drie aflezingen bij $\bar y$:
$0.0838\sqrt{1/3 + 1/7} = \qty{0.058}{mM}$.
\end{solution}

\begin{solution}{pb:b3:electrode-kinetics:1}
\textbf{1.} $\ce{C6H12O6 -> C6H10O6 + 2 H+ + 2 e-}$ en $\ce{[Fe(CN)6]^3- + e- -> [Fe(CN)6]^4-}$;
totaal
\[
  \ce{C6H12O6 + 2 [Fe(CN)6]^3- -> C6H10O6 + 2 [Fe(CN)6]^4- + 2 H+} .
\]
Aan de elektrode: $\ce{[Fe(CN)6]^4- -> [Fe(CN)6]^3- + e-}$.
\textbf{2.} Twee.
\textbf{3.} De opgeloste zuurstof in bloed is laag en variabel, en haar
product, waterstofperoxide, wordt pas geoxideerd bij een hoge potentiaal
waar veel andere stoffen reageren; de mediator is in een bekende overmaat
aanwezig.
\textbf{4.} Zodat zijn oppervlakteconcentratie nul is: de stroom wordt dan
alleen door diffusie beperkt en is ongevoelig voor kleine
potentiaalveranderingen.
\textbf{5.} De stroom daalt als $t^{-1/2}$; de aflezingen moeten gebeuren op
het tijdstip dat voor de ijking gebruikt werd.
\textbf{6.} Cottrell: $i\sqrt t = 42.0$, 41.4, 41.7, 41.6, 41.8: constant
binnen 1\,\%.
\textbf{7.} $D = \pi(\text{helling}/FAc)^2 = \pi(\num{41.8e-6}/(96\,485 \times 0.0300 \times \num{1.00e-5}))^2 = \qty{6.55e-6}{cm^2.s^{-1}}$.
\textbf{8.} $\sqrt{\pi \times \num{6.55e-6} \times 5} = \qty{0.010}{cm} = \qty{100}{\micro m}$:
de verarming bereikt de bovenkant van de laag na ongeveer \qty{5}{s};
latere aflezingen zouden onder de wet van Cottrell vallen.
\textbf{9.} De helft: \qty{9.35}{\micro A}.
\textbf{10.} De achtergrondstroom: oxidatie van storende stoffen en
onzuiverheden, opladen, het restje hexacyanoferraat(II) in de mediator
zelf.
\textbf{11.} Ja: de residuen, hoogstens \qty{0.095}{\micro A}, spreiden
zonder trend. Bij veel glucose raakt de mediator, die in een beperkte
hoeveelheid aanwezig is, of het \omterm{def:b3:complex-kinetics:enzyme}{enzym} verzadigd en buigt de rechte af.
\textbf{12.} $(7.10 - 0.356)/0.9189 = \qty{7.34}{mM}$.
\textbf{13.} $0.0838\sqrt{1 + 1/7 + (7.10 - 8.89)^2/203.5} = 0.0838 \times 1.076 = \qty{0.090}{mM}$.
\textbf{14.} De term $1/m$ (één enkele aflezing); herhaalde aflezingen, of
meerdere strookjes, verkleinen hem.
\textbf{15.} $3 \times 0.077/0.919 = \qty{0.25}{mM}$.
\textbf{16.} $7.34 \times 180.16 = \qty{1322}{mg.L^{-1}} = \qty{132}{mg.dL^{-1}}$.
\textbf{17.} Het voegt zijn eigen oxidatiestroom toe: de meter leest te hoog
af.
\textbf{18.} Een \omterm{def:b3:electrode-kinetics:cv}{voltammogram} van het strookje zonder glucose toont een
anodische golf van ascorbaat waar hexacyanoferraat(II) geoxideerd wordt.
\textbf{19.} Bij een lagere potentiaal worden minder stoffen geoxideerd.
\textbf{20.} De blancostroom wordt van de aflezing van de enzymelektrode
afgetrokken voordat de ijklijn toegepast wordt.
\textbf{21.} $D$ groeit met enkele procenten per kelvin en de stroom als
$\sqrt D$: meters meten de temperatuur en corrigeren ervoor.
\textbf{22.} De viscositeit en het aandeel rode bloedcellen van het bloed
veranderen $D$ en de stroom; standaarden in een gelijkaardige matrix heffen
deze effecten op.
\textbf{23.} Variaties van strookje tot strookje (elektrodeoppervlak,
belading met \omterm{def:b3:complex-kinetics:enzyme}{enzym} en mediator), temperatuur en samenstelling van het
bloed komen bij de eigen onzekerheid van de ijking.
\textbf{24.} \textbf{$c(\text{glucose}) = 7.34 \pm \qty{0.09}{mmol.L^{-1}}$}
(standaardonzekerheid), ongeveer \qty{132}{mg.dL^{-1}}.
\end{solution}
